

\section{Introduction}
\label{sec:intro}
Novel view synthesis (NVS) stands as a critical task in the domain of 3D vision and plays a vital role in many applications, \eg VR, AR, and movie production. NVS aims at rendering images from any desired viewpoint or timestamp of a scene, usually requiring modeling the scene accurately from several 2D images. Dynamic scenes are quite common in real scenarios, rendering which is important but challenging as complex motions need to be modeled with both spatially and temporally sparse input.

NeRF~\cite{mildenhall2021nerf} has achieved great success in synthesizing novel view images by representing scenes with implicit functions. The volume rendering techniques~\cite{drebin1988volume} are introduced to connect 2D images and 3D scenes. However, the original NeRF method bears big training and rendering costs. Though some NeRF variants~\cite{dvgo,tineuvox,hexplane,kplanes,tensor4d,fridovich2022plenoxels,instantngp} reduce the training time from days to minutes, the rendering process still bears a non-negligible latency.

Recent 3D Gaussian Splatting (3D-GS)~\cite{3dgs} significantly boosts the rendering speed to a real-time level by representing the scene as 3D Gaussians. The cumbersome volume rendering in the original NeRF is replaced with efficient differentiable splatting~\cite{yifan2019differentiablesplatting}, which directly projects 3D Gaussian points onto the 2D plane. 3D-GS not only enjoys real-time rendering speed but also represents the scene more explicitly, making it easier to manipulate the scene representation. 

However, 3D-GS focuses on the static scenes. Extending it to dynamic scenes as a 4D representation is a reasonable, important but difficult topic. The key challenge lies in modeling complicated point motions from sparse input. 3D-GS holds a natural geometry prior by representing scenes with point-like Gaussians. One direct and effective extension approach is to construct 3D Gaussians at each timestamp~\cite{dynamic3dgs} but the storage/memory cost will multiply especially for long input sequences. Our goal is to construct a compact representation while maintaining both training and rendering efficiency, \ie 4D Gaussian Splatting (\textbf{4D-GS}). To this end, we propose to represent Gaussian motions and shape changes by an efficient Gaussian deformation field network, containing a temporal-spatial structure encoder and an extremely tiny multi-head Gaussian deformation decoder. Only one set of canonical 3D Gaussians is maintained. For each timestamp, the canonical 3D Gaussians will be transformed by the Gaussian deformation field into new positions with new shapes. The transformation process represents both the Gaussian motion and deformation. Note that different from modeling motions of each Gaussian separately~\cite{dynamic3dgs,4dgs-2}, the spatial-temporal structure encoder can connect different adjacent 3D Gaussians to predict more accurate motions and shape deformation. Then the deformed 3D Gaussians can be directly splatted for rendering the according-timestamp image. Our contributions can be summarized as follows.

\begin{itemize}
     \item An efficient 4D Gaussian splatting framework with an efficient Gaussian deformation field is proposed by modeling both Gaussian motion and Gaussian shape changes across time.
     \item A multi-resolution encoding method is proposed to connect the nearby 3D Gaussians and build rich 3D Gaussian features by an efficient spatial-temporal structure encoder. 
     \item 4D-GS achieves real-time rendering on dynamic scenes, up to 82 FPS at a resolution of 800$\times$800 for synthetic datasets and 30 FPS at a resolution of 1352$\times$1014 in real datasets, while maintaining comparable or superior performance than previous state-of-the-art (SOTA) methods and shows potential for editing and tracking in 4D scenes.
     % This performance surpasses that of other implementations.
\end{itemize}